\documentclass[11pt,a4paper]{article}

\usepackage[margin=1in]{geometry}
\usepackage{amsmath,amssymb,amsthm,bm}
\usepackage{booktabs}
\usepackage{graphicx}
\usepackage{xcolor}
\usepackage{listings}
\usepackage[protrusion=true,expansion=false]{microtype}
\usepackage{enumitem}
\usepackage{caption}
\usepackage[colorlinks=true,linkcolor=darkblue,citecolor=darkblue,urlcolor=darkblue]{hyperref}

\definecolor{darkblue}{RGB}{20,60,120}
\definecolor{codebg}{RGB}{248,248,248}
\definecolor{rulegray}{RGB}{200,200,200}

\lstset{
  basicstyle=\ttfamily\footnotesize,
  backgroundcolor=\color{codebg},
  breaklines=true,
  frame=single,
  rulecolor=\color{rulegray},
  framesep=5pt,
  language=Python,
  keywordstyle=\color{darkblue}\bfseries,
  commentstyle=\color{gray}\itshape,
  stringstyle=\color{red!50!black},
  showstringspaces=false,
  captionpos=b,
}

\newtheorem{definition}{Definition}
\newtheorem{proposition}{Proposition}
\newtheorem{remark}{Remark}

\title{\vspace{-2em}\textbf{Bayesian Orchestration of Tool-Augmented Agents}\\[0.4em]
\large Sequential Monte Carlo over Action Trajectories,\\ with Application to Retrieval-Augmented Generation}

\author{Angshul Majumdar\\
\small Department of Electronics and Communication Engineering\\
\small Indraprastha Institute of Information Technology Delhi\\
\small \texttt{angshul@iiitd.ac.in}}

\date{\today}

\begin{document}
\maketitle

\begin{abstract}
\noindent
Tool-augmented language model agents select actions greedily and recover from
failure by retrying. Both behaviours follow from a single design decision: the
agent maintains exactly one execution trajectory and discards the alternatives.
Evidence that arrives later---a failed validation, a contradicting
source---then has no surviving branch to promote, and a retry redraws from the
same proposal distribution rather than conditioning on what was observed. It is
a draw from the prior, not a posterior update.

We formulate agent execution as sequential inference over action trajectories
under a finite-horizon generative model, and show that three familiar
behaviours---deterministic greedy selection, probabilistic forward sampling,
and Bayesian sequential Monte Carlo (SMC)---are inference strategies over
\emph{one} model rather than distinct architectures. Making the inference
regime a runtime parameter permits controlled comparison: identical tools,
proposals, and likelihoods run under all three, so any difference in outcome is
attributable to inference alone.

Each particle additionally carries private Beta--Bernoulli beliefs over tool
reliability, updated online with \emph{retroactive credit assignment}: a
checker's verdict at step $t$ revises the standing of whichever tool produced
the claim it examined at step $t-1$. We further show that answer selection must
be governed by validation status rather than recency, and that pooling
posterior mass by answer rather than by trajectory yields a calibration signal
that identifies when the evidence does not settle a question---at which point
the agent asks rather than guesses.

Empirically, on stochastically unreliable tools the method converts a source
correct one call in ten into a system correct nine times in ten. On three
retrieval-augmented generation failure modes---coreference, staleness under
vocabulary drift, and index noise---evaluated against BM25, TF-IDF/cosine, and
LSA, it achieves $1.000$, $0.903$, and $1.000$ against a conventional
single-pass pipeline's two deterministic failures and $0.420$. We also report
where the method does \emph{not} help: a correctly engineered conditional graph
is competitive and cheaper when the failure mode is known in advance, and a
confidently misspecified prior costs more than it costs a system holding no
prior at all.

\vspace{0.5em}
\noindent\textbf{Software.} MIT-licensed, no required runtime dependencies,
$218$ tests. \url{https://github.com/AngshulMajumdar/BayesianRAG}
\end{abstract}

\vspace{0.5em}
\hrule
\vspace{1em}

\section{Introduction}

An agent built on a language model interleaves generation with calls to
external tools: retrievers, calculators, database queries, verification
services. Execution is therefore a sequence of decisions under uncertainty in
which the consequence of a choice is frequently unobservable until several
steps later.

Mainstream frameworks do not treat it this way. Action selection is implemented
as deterministic parsing of model output followed by a greedy choice, and
failure is handled by retrying. These are the same decision seen twice: the
agent maintains one trajectory. When a verification step contradicts an earlier
retrieval, no alternative survives for that evidence to promote. Retrying does
not repair this, because a retry samples afresh from the same proposal without
conditioning on the intervening observation.

The consequences are ordinary and familiar. An agent commits to a fast,
plausible, stale source and never revisits it. A retrieval pipeline returns the
top-ranked passage when the answer sits at rank two. A pipeline that
\emph{does} validate discards the validated answer because it later looked
somewhere else. Each is a failure of \emph{inference}, not of the underlying
tools.

\subsection{Contributions}

\begin{enumerate}[leftmargin=1.4em,itemsep=0.25em]
\item A formulation of agent execution as sequential inference over action
      trajectories, under which greedy selection, forward sampling, and SMC
      are inference strategies over one shared generative model
      (\S\ref{sec:model}--\S\ref{sec:ladder}).
\item A per-particle Beta--Bernoulli reliability model with retroactive credit
      assignment, letting a late verdict revise the standing of an earlier
      decision (\S\ref{sec:reliability}).
\item An answer-selection rule ordered by validation status rather than
      recency, and a demonstration that the natural alternative silently
      discards confirmed answers (\S\ref{sec:selection}).
\item Marginalisation of posterior mass over \emph{answers} rather than
      trajectories, yielding a consensus statistic and a principled criterion
      for requesting clarification (\S\ref{sec:consensus}).
\item Controlled experiments isolating the contribution of inference, including
      results that do not favour the method (\S\ref{sec:experiments}).
\end{enumerate}

\section{Agent Execution as Sequential Inference}
\label{sec:model}

\subsection{Generative model}

Let $a_t \in \mathcal{A}$ denote the action taken at step $t$ (a tool
invocation with arguments) and $o_t \in \mathcal{O}$ the resulting observation.
Write the history $h_t = (a_{1:t-1}, o_{1:t-1})$, with $h_1 = \emptyset$. An
episode of horizon $T$ has joint density
\begin{equation}
p(a_{1:T}, o_{1:T})
  \;=\; \prod_{t=1}^{T} \underbrace{p(a_t \mid h_t)}_{\text{proposal}}\;
                        \underbrace{p(o_t \mid a_t, h_t)}_{\text{likelihood}}.
\label{eq:joint}
\end{equation}

The proposal $p(a_t \mid h_t)$ is supplied by the model or policy: it is what a
language model expresses when it emits a tool call. The likelihood
$p(o_t \mid a_t, h_t)$ scores how well an observation coheres with the
trajectory that produced it---high when a checker confirms a claim, low when it
refutes one.

The quantity a practitioner wants is the posterior over trajectories given what
the tools actually returned,
\begin{equation}
p(a_{1:T} \mid o_{1:T}) \;\propto\; \prod_{t=1}^{T} p(a_t \mid h_t)\, p(o_t \mid a_t, h_t).
\label{eq:posterior}
\end{equation}

\begin{remark}
Equation~\eqref{eq:joint} is deliberately weak. It assumes only that actions
are proposed conditional on history and that observations are scored
conditional on the action and history. It does not assume Markov structure,
a fixed tool set, or a differentiable policy.
\end{remark}

\subsection{Why greedy execution cannot recover}

A conventional agent computes
$\hat a_t = \arg\max_{a} p(a \mid h_t)$
and proceeds. This is maximum \emph{a priori} selection: the likelihood
$p(o_t \mid a_t, h_t)$ is never consulted. Because a single trajectory is
retained, the posterior~\eqref{eq:posterior} is approximated by a degenerate
measure at $\hat a_{1:T}$, and no observation can shift mass onto a branch that
was never instantiated.

\begin{proposition}[Greedy irrecoverability]
Let $\hat a_{1:T}$ be the greedy trajectory. For any observation sequence
$o_{1:T}$, the greedy estimate is independent of $o_{1:T}$ except through its
effect on the proposal $p(a_t\mid h_t)$. In particular, if the proposal at
step $t$ does not depend on $o_{t-1}$, then no evidence obtained at step $t-1$
can alter the action taken at step $t$.
\end{proposition}

This is why retrying does not help. A retry re-evaluates
$p(a_t \mid h_t)$ with $h_t$ largely unchanged and yields the same action, or a
sample from the same distribution. It is a draw from the prior.

\section{The Inference Ladder}
\label{sec:ladder}

Three behaviours arise as inference strategies over
Equation~\eqref{eq:joint}. They differ in two binary properties: whether
multiple hypotheses are instantiated (\emph{explore}), and whether the
observation likelihood is accumulated into their weights (\emph{reweight}).

\begin{table}[h]
\centering
\begin{tabular}{@{}llccl@{}}
\toprule
\textbf{Regime} & \textbf{Particles} & \textbf{Explore} & \textbf{Reweight} & \textbf{Approximates} \\
\midrule
Deterministic greedy   & $1$ & no  & no  & point mass at the modal prior path \\
Probabilistic forward  & $N$ & yes & no  & samples from the \emph{prior} $p(a_{1:T})$ \\
Bayesian SMC           & $N$ & yes & yes & the posterior $p(a_{1:T}\mid o_{1:T})$ \\
\bottomrule
\end{tabular}
\caption{The inference ladder. Only SMC targets the posterior.}
\label{tab:ladder}
\end{table}

The distinction between the second and third rows is the one most easily
conflated in practice: \emph{maintaining several hypotheses is not the same as
performing inference over them}. Forward sampling explores, but because weights
are never revised, its particle cloud remains a sample from the prior. \S\ref{sec:exp-ladder}
shows this gap is large empirically---roughly a coin flip versus recovery in
nine runs out of ten on the same tools.

\subsection{Sequential Monte Carlo over trajectories}

We approximate the posterior by a weighted particle set
$\{(a_{1:t}^{(i)}, w_t^{(i)})\}_{i=1}^{N}$ with
$\hat p(a_{1:t} \mid o_{1:t}) = \sum_i \bar w_t^{(i)} \delta_{a_{1:t}^{(i)}}$,
where $\bar w_t^{(i)}$ are the normalised weights. At each step, particle $i$
samples $a_t^{(i)} \sim q(\cdot \mid h_t^{(i)})$, executes it, and updates its
weight by the incremental factor
\begin{equation}
w_t^{(i)} \;=\; w_{t-1}^{(i)} \cdot
  \frac{p(o_t^{(i)} \mid a_t^{(i)}, h_t^{(i)})\, p(a_t^{(i)} \mid h_t^{(i)})}
       {q(a_t^{(i)} \mid h_t^{(i)})}.
\label{eq:weight}
\end{equation}
Taking the proposal $q$ equal to the model's own action distribution
$p(a_t \mid h_t)$---the natural choice, since that is what the policy
emits---the ratio collapses and Equation~\eqref{eq:weight} reduces to
\begin{equation}
\log w_t^{(i)} \;=\; \log w_{t-1}^{(i)} + \log p(o_t^{(i)} \mid a_t^{(i)}, h_t^{(i)}),
\label{eq:logweight}
\end{equation}
accumulated in log space for numerical stability. All arithmetic uses the
max-subtraction identity
$\log \sum_i e^{x_i} = x^\ast + \log \sum_i e^{x_i - x^\ast}$,
$x^\ast = \max_i x_i$, so weights spanning hundreds of orders of magnitude
normalise without overflow.

\subsection{Degeneracy and resampling}

Particle weights concentrate over time. We monitor the effective sample size
\begin{equation}
\mathrm{ESS}_t \;=\; \Big( \textstyle\sum_{i=1}^{N} (\bar w_t^{(i)})^2 \Big)^{-1}
\;\in\; [1, N],
\end{equation}
and resample when $\mathrm{ESS}_t / N$ falls below a threshold $\tau$
(default $0.4$), using systematic resampling---which has lower variance than
multinomial resampling and $O(N)$ cost---after which weights are reset to
uniform. Resampling is suppressed on the final step, where it would add
variance without informing any subsequent decision.

\subsection{The diversity requirement}

\begin{proposition}[Selector degeneracy]
\label{prop:diversity}
If the selection rule is deterministic given a particle's state, and all
particles share an identical initial state, then all $N$ particles follow
identical trajectories. The weighted particle set is then a point mass, and
reweighting has no effect: SMC reduces exactly to greedy execution at $N$
times the cost.
\end{proposition}

This is not a hypothetical. An $\arg\max$ selector produces a cloud of
identical particles, and the resulting system silently behaves as a slow greedy
agent while appearing to run a particle filter. Selection must therefore sample
in proportion to its score under exploring regimes. In the reference
implementation the regime is passed to the selector precisely so this cannot be
got wrong by accident, and a regression test asserts that greedy and SMC
produce materially different trajectory distributions.

\section{Tool Reliability and Retroactive Credit}
\label{sec:reliability}

\subsection{Conjugate belief}

Each tool $j$ has an unknown success probability $\theta_j$ with a conjugate
Beta prior. After $k$ successes in $n$ trials,
\begin{equation}
p(\theta_j \mid \mathcal{D}) = \mathrm{Beta}(\alpha_j + k,\; \beta_j + n - k),
\qquad
\mathbb{E}[\theta_j \mid \mathcal{D}] = \frac{\alpha_j + k}{\alpha_j + \beta_j + n}.
\end{equation}
Beliefs are held \emph{per particle}, not globally. A trajectory burned by an
unreliable tool distrusts it, while a sibling that never invoked it does not---
which is the correct semantics, since the two have observed different evidence.

Selection scores a candidate by the product of its posterior mean and a static
\emph{appeal} $\rho_j$ encoding cost, latency, or convenience:
\begin{equation}
s_j \;=\; \mathbb{E}[\theta_j \mid \mathcal{D}] \cdot \rho_j .
\label{eq:score}
\end{equation}
Separating $\rho_j$ from $\theta_j$ is deliberate. The failure mode this work
addresses is precisely a tool that is \emph{attractive and wrong}: high $\rho$,
low $\theta$. Collapsing them into one number makes that case
inexpressible.

\subsection{Retroactive credit assignment}

The mechanism by which late evidence revises an earlier decision is that a
scorer may emit belief updates against tools \emph{other than} the one just
executed. When a checker examines a claim and rejects it, the update is applied
to whichever tool produced that claim:
\begin{equation}
\theta_{j^\ast} \leftarrow \mathrm{Beta}(\alpha_{j^\ast},\, \beta_{j^\ast} + \lambda),
\qquad
j^\ast = \text{the tool that produced the examined claim},
\end{equation}
with $\lambda$ a credit weight (default $2$). A source is thus penalised when
it is \emph{caught}, not when it answers.

\begin{remark}[Per-particle claim resolution]
A subtlety with real consequences: the proposer runs once per step and emits
one candidate menu shared by all particles, but each particle must check
\emph{its own} claim. Binding the checker's argument to any single particle's
claim causes cross-particle misattribution, under which a particle using a
reliable tool is penalised for a claim produced elsewhere in the cloud---
driving a correct tool's posterior monotonically downward while it continues to
be correct. The implementation resolves a \texttt{Pending.CLAIM} placeholder
per particle at execution time.
\end{remark}

\subsection{Cross-episode transport}

Beliefs may be carried between episodes with a tempering factor
$\gamma \in [0,1]$ applied to accumulated evidence:
\begin{equation}
(\alpha', \beta') \;=\; \big(\alpha_0 + \gamma(\alpha - \alpha_0),\;
                             \beta_0 + \gamma(\beta - \beta_0)\big).
\end{equation}
This transports the posterior itself. Reconstructing a Beta from its mean
alone---the obvious shortcut---preserves the estimate while destroying the
sample size behind it, converting fifty confident observations into a
near-uninformative prior. Note that $\gamma < 1$ also shrinks the mean toward
the prior, which is correct rather than incidental: with less evidence the
estimate should be both less certain and closer to where it began.

\section{Answer Selection Under Validation}
\label{sec:selection}

Given a terminal particle, which of its observations constitutes the answer?
The obvious rule---take the most recent---is wrong in a way that is easy to
miss and expensive in practice.

Let $\mathcal{V}$ denote the set of claims a checker in this trajectory
confirmed, and $\mathcal{R}$ the set it rejected. The correct preference order
is
\begin{equation}
\text{answer} \;=\;
\begin{cases}
a \in \mathcal{V} & \text{if } \mathcal{V} \neq \emptyset,\\[0.3em]
\text{most recent } a \notin \mathcal{R} & \text{else if such } a \text{ exists},\\[0.3em]
\text{most recent } a & \text{otherwise.}
\end{cases}
\label{eq:answer}
\end{equation}

Both refinements matter. Omitting the $\mathcal{R}$ clause lets an agent return
a claim its own checker disproved---it gathered decisive evidence and then
ignored it. Omitting the $\mathcal{V}$ clause is subtler: a trajectory that
found the correct passage, validated it, and subsequently examined something
else will return whatever it touched last, discarding a confirmation it had
already paid for. Empirically this second error was worth between $3$ and $35$
percentage points depending on the task, and---because adding particles raises
the chance that \emph{some} trajectory happens to terminate on the validated
answer---it masquerades convincingly as a budget limitation.

\section{Posterior Consensus and Clarification}
\label{sec:consensus}

Trajectory mass is the wrong measure of agreement. Sixteen particles that reach
one answer by sixteen different routes agree completely, yet each holds
$1/16$ of the mass. We therefore marginalise over answers,
\begin{equation}
\pi(a) \;=\; \sum_{i=1}^{N} \bar w_T^{(i)} \,
             \mathbb{1}\!\left[\mathrm{answer}(a_{1:T}^{(i)}) = a\right],
\qquad
\mathrm{consensus} \;=\; \max_a \pi(a),
\end{equation}
which distinguishes genuine disagreement from mere path diversity. When
\begin{equation}
\mathrm{consensus} < \kappa \quad\text{and}\quad |\{a : \pi(a) > 0\}| > 1
\end{equation}
for a threshold $\kappa$ (default $0.6$), no answer commands enough of the
posterior to be asserted, and the agent requests clarification rather than
selecting arbitrarily among incompatible answers.

This is the state a single-path pipeline cannot represent. Having committed to
one trajectory it holds exactly one answer and has no representation for the
proposition that the evidence supports several. In our experiments the
criterion fires in $13$ of $200$ runs where a checker settles the question, and
in essentially all runs where two sources irreconcilably disagree---it asks
when genuinely split, not by default.

\section{Implementation}

The reference implementation separates four concerns. A \emph{proposer} emits
candidate actions; a \emph{selector} chooses among them per particle; an
\emph{environment} executes; a \emph{scorer} maps observations to likelihoods
and belief updates. A single runner implements
\[
\text{propose} \rightarrow \text{select} \rightarrow \text{observe}
\rightarrow \text{weight} \rightarrow \text{resample if } \mathrm{ESS}\ \text{collapses}
\]
for every regime; the regime supplies only a particle count and the two boolean
flags of Table~\ref{tab:ladder}.

For the common case all three components are derived from tool metadata, so a
working agent requires only tool declarations:

\begin{lstlisting}[caption={The complete high-level API.}]
from bayesian_rag import Agent, tool, checker

@tool(reliability=0.6, appeal=0.9)     # attractive but unreliable
def quick_search(query: str) -> str:
    """Fast, sometimes stale."""
    ...

@tool(reliability=0.95, appeal=0.5)    # unglamorous but correct
def verified_search(query: str) -> str:
    """Slow, authoritative."""
    ...

@checker
def fact_check(text: str) -> bool:
    """Validate a claim."""
    ...

agent  = Agent([quick_search, verified_search, fact_check])
result = agent.run("What is the capital of Australia?")

result.answer                # maximum a posteriori answer
result.consensus             # posterior mass backing it
result.status                # validated | unverified | refuted
result.needs_clarification   # evidence insufficient to assert
result.reliability           # learned posterior means per tool
\end{lstlisting}

\paragraph{Determinism and memoisation.}
Identical invocations are memoised within an episode, so $N$ particles over a
$k$-tool menu issue $k$ calls per step rather than $N k$. This is sound only
for deterministic tools. A stochastic tool must be declared
\texttt{deterministic=False}; otherwise every particle receives the same cached
draw, the cloud collapses to a single sample, and the re-drawing mechanism that
lets a filter exceed an unreliable tool's accuracy is silently disabled. This
was a real defect during development and is now covered by regression tests.

\paragraph{Reproducibility.}
Runs are exactly reproducible from a seed, and \emph{only} from the seed: a
test asserts that reseeding the global generator cannot perturb a seeded run,
since a component reaching for global randomness breaks replication in a way
that surfaces only when someone attempts to reproduce a result.

\paragraph{Statistics.}
All reported proportions carry Wilson score intervals, which remain within
$[0,1]$ and behave correctly near the boundaries---necessary here, since
several conditions score exactly $0.000$ or $1.000$ where the normal
approximation returns impossible bounds. Differences use Newcombe intervals and
a two-proportion $z$-test. All are implemented in the standard library; the
core package has no required runtime dependencies.

\section{Experiments}
\label{sec:experiments}

All experiments use deterministic seeds indexed by trial, so every figure below
regenerates exactly. Tools are synthetic, with i.i.d. Bernoulli success---an
assumption we return to in \S\ref{sec:limitations}.

\subsection{The inference ladder}
\label{sec:exp-ladder}

Two adversarial scenarios pit a cheap, confident, incorrect source against a
slower authoritative one, with a checker that reveals the conflict one step
later. A third scenario, with unambiguous routing, acts as a control. Crucially
the proposer offers the same menu at every step: no recovery is scripted into
the control flow, so any recovery observed is produced by inference.

\begin{table}[h]
\centering
\begin{tabular}{@{}lccc@{}}
\toprule
\textbf{Scenario} & \textbf{Greedy} & \textbf{Forward} & \textbf{SMC} \\
\midrule
\texttt{stale\_vs\_verified}          & $0.000$ & $0.510$ & $\bm{0.957}$ \\
\texttt{web\_vs\_official}            & $0.000$ & $0.520$ & $\bm{0.947}$ \\
\texttt{ambiguous\_location} (control)& $1.000$ & $1.000$ & $1.000$ \\
\bottomrule
\end{tabular}
\caption{Success rate, $300$ seeds per cell. All contrasts on the adversarial
scenarios are significant at $p < 10^{-10}$; none on the control are
significant.}
\end{table}

Greedy fails \emph{completely}, not merely often: with one particle and an
$\arg\max$ selector it re-selects the attractive wrong source at every step,
and no evidence can dislodge it. The forward-versus-SMC gap is the substantive
result---both explore identically and differ only in whether the likelihood
reaches the weights, and that alone converts a coin flip into recovery. The
control confirms the machinery costs nothing where routing is already
unambiguous.

\subsection{Unreliable and adversarial tools}

The regime where inference contributes most is one no fixed routing rule can
address: tools that are correct \emph{sometimes}, with no stable pattern. There
is no correct branch to hard-code, because which source is right varies call to
call.

\begin{table}[h]
\centering
\begin{tabular}{@{}lccc@{}}
\toprule
\textbf{Tool accuracies} & \textbf{LangGraph (engineered)} & \textbf{LangChain (reflexive)} & \textbf{SMC} \\
\midrule
one at $40\%$                 & $0.607$ & $0.407$ & $\bm{0.933}$ \\
two at $40\%$                 & $0.607$ & $0.607$ & $\bm{0.887}$ \\
$30\% + 90\%$                 & $0.913$ & $0.913$ & $\bm{0.960}$ \\
$10\% + 40\% + 85\%$          & $0.820$ & $0.820$ & $\bm{0.867}$ \\
three at $35\%$               & $0.567$ & $0.567$ & $\bm{0.787}$ \\
\bottomrule
\end{tabular}
\caption{Stochastically unreliable tools, $150$ seeds per cell.}
\end{table}

The mechanism is independent re-drawing: each particle samples afresh, a
checker validates, and the filter promotes whichever draw survives. It scales
sharply---at $16$ particles and $6$ steps, a source correct $10\%$ of the time
yields a system correct $93.5\%$ of the time; $64$ particles convert a $35\%$
source to $99\%$. A greedy agent cannot do this: one trajectory obtains one
draw, with nothing to compare it against.

\subsection{Prior sensitivity}

The method inherits whatever reliability the practitioner declares. Against a
tool with true accuracy $0.20$:

\begin{center}
\begin{tabular}{@{}lccccc@{}}
\toprule
Declared prior & $0.2$ (honest) & $0.5$ & $0.9$ & $0.99$ (confident lie) \\
\midrule
SMC accuracy & $\bm{0.993}$ & $0.887$ & $0.527$ & $0.380$ \\
\bottomrule
\end{tabular}
\end{center}

A $61$-point swing from the declaration alone. Inspection of the agent's seeded
selection scores identifies the mechanism: at a high declared prior the tool's
score~\eqref{eq:score} exceeds the checker's, so the agent re-invokes the
unreliable source \emph{instead of validating it}. A confidently false prior
does not merely misroute; it makes the agent stop checking its own work.

\subsection{Comparison with deterministic orchestration}

Results here do not uniformly favour the method, and we report them as
measured. A naive pipeline is not an interesting baseline---it has no recovery
path at all. The honest comparison is against a conditional graph written by
someone who \emph{knows} the failure mode: validate, and escalate on failure.

Across three worlds sharing one tool set---(A) the anticipated failure, (B) the
same tools with roles reversed, (C) the escalation target unavailable---the
engineered graph passes all three, at lower cost, deterministically. It passes
(C) by crashing rather than reporting uncertainty, and its success in (B) is
luck: its hard-coded preference happens to match reality. But it succeeds.

SMC with \emph{honest} priors also passes all three; SMC with the declared
(and, in world B, false) priors scores $0.31$ and loses to every deterministic
baseline. The conclusion we draw is narrow and worth stating plainly: use a
conditional edge when the failure mode can be named in advance; reach for
inference when it cannot---and only if the priors are honest, because the
posterior faithfully propagates a confident lie.

Tool-call count is \emph{not} reported as a finding. Maintaining a belief over
alternatives costs more calls by construction; the question is what they buy.

\subsection{Retrieval-augmented generation}

We evaluate three RAG failure modes against three retrievers people actually
deploy: BM25 (Okapi, $k_1=1.5$, $b=0.75$), TF-IDF with cosine similarity, and
LSA (TF-IDF with truncated SVD). Every instance is \emph{verified} to defeat
every retriever before its result is reported.

\begin{table}[h]
\centering
\small
\begin{tabular}{@{}lcccc@{}}
\toprule
\textbf{Instance} & \textbf{Normal RAG} & \textbf{Bayesian (BM25)} & \textbf{(TF-IDF)} & \textbf{(LSA)} \\
\midrule
1. Coreference             & wrong, always & $\bm{1.000}$ & $\bm{1.000}$ & $0.967$ \\
2. Staleness (vocab drift) & wrong, always & $\bm{0.903}$ & $\bm{0.903}$ & $\bm{0.903}$ \\
3. Noisy index             & $0.420$       & $\bm{1.000}$ & $\bm{1.000}$ & $\bm{1.000}$ \\
\bottomrule
\end{tabular}
\caption{$300$ seeds, $95\%$ Wilson intervals. Normal RAG is a single
retrieval call, top-1, no checking.}
\end{table}

\textbf{Coreference.} The answer-bearing passage refers to the entity as
``the company'' and never repeats its name, while a non-answering passage names
it three times. BM25 scores them $3.06$ versus $0.93$; the answer lands at rank
two. No lexical retriever performs coreference resolution. This is the most
common real RAG failure---\emph{the answer is in top-$k$ but not top-1}---and we
verify it is not a small-corpus artifact: the failure holds from $3$ to $52$
documents with the score gap widening.

\textbf{Staleness under vocabulary drift.} An outdated passage uses the
canonical phrasing the query mirrors (``generating capacity''); the current
amendment uses internal shorthand (``output rating''). BM25 scores them $5.40$
versus $0.22$. Lexical scoring has no notion of recency.

\textbf{Index noise.} The retriever is degraded to model replication lag.
Normal RAG's $0.420$ matches the closed form
$0.9 \times \tfrac{1}{3} + 0.1 \times 1.0 = 0.40$. The Bayesian pipeline
declares the retrieval tool non-deterministic, so each particle draws
independently---as when re-querying a replicated cluster.

\section{Limitations}
\label{sec:limitations}

We state these plainly, and several arose from findings that contradicted our
own earlier claims.

\begin{itemize}[leftmargin=1.4em,itemsep=0.3em]
\item \textbf{Synthetic tools.} All tools are mocks with i.i.d. Bernoulli
      success. Real failures correlate with query type, load, and upstream
      incidents in ways a fixed Bernoulli does not capture.
\item \textbf{Perfect checkers.} Verifiers here compare against ground truth
      exactly. A real verifier is itself unreliable, which would compress every
      gap reported above. Two genuine checker bugs were found during this work;
      both are documented in the repository rather than quietly fixed.
\item \textbf{Comparison baselines are semantic stand-ins}, not the real
      LangChain and LangGraph packages, which could not be installed in the
      evaluation environment. The engineered graph's agreement with its
      closed-form expectation is partial evidence its control flow is correct.
      An adapter is provided to verify against the real packages.
\item \textbf{Scenarios were constructed to separate the regimes.} Recovery is
      not scripted into the control flow, but the scenarios were designed to
      exhibit the failure modes of interest. They establish that the mechanism
      works where it matters, not how often that situation arises.
\item \textbf{Priors are inherited, not learned from scratch.} A confidently
      wrong prior is costly, as \S\ref{sec:experiments} quantifies.
\item \textbf{Short horizons and small corpora.} Episodes run to a few steps;
      corpora contain a handful of documents. Particle degeneracy over long
      horizons is not stressed, and BM25's IDF statistics behave differently at
      scale.
\item \textbf{Dense retrieval is only partially tested.} LSA falsified our own
      earlier conjecture that dense retrieval would resolve the coreference
      instance---it ranks the answer \emph{worse} than BM25. Whether a
      pretrained sentence transformer would resolve it remains open.
\end{itemize}

\section{Conclusion}

Treating agent orchestration as inference over trajectories, rather than as
control flow with retries, makes a specific and testable difference: evidence
obtained late in an episode can revise a decision made early in it. Formulating
greedy execution, forward sampling, and SMC as strategies over one generative
model makes the comparison controlled, since the regime is the only thing that
varies.

The empirical picture is not uniform, and we regard that as the more useful
result. Where a failure mode can be named in advance, a conditional edge is
cheaper, deterministic, and sufficient. Where it cannot---unreliable sources,
conflicting evidence, retrieval whose ranking is wrong in ways not known
beforehand---maintaining a posterior over alternatives converts sources that
are individually unreliable into systems that are collectively dependable, and
supplies a calibrated signal for when to stop guessing and ask.

\section*{Availability}

Software, experiments, figures, and raw results:
\url{https://github.com/AngshulMajumdar/BayesianRAG}. MIT licensed. The core
package has no required runtime dependencies. Every table in this paper is
regenerated by a documented command.

\begin{thebibliography}{99}
\small

\bibitem{doucet2009} A.~Doucet and A.~M. Johansen.
A tutorial on particle filtering and smoothing: fifteen years later.
\emph{Handbook of Nonlinear Filtering}, Oxford University Press, 2009.

\bibitem{chopin2020} N.~Chopin and O.~Papaspiliopoulos.
\emph{An Introduction to Sequential Monte Carlo}. Springer, 2020.

\bibitem{delmoral2006} P.~Del~Moral, A.~Doucet, and A.~Jasra.
Sequential Monte Carlo samplers.
\emph{Journal of the Royal Statistical Society: Series B}, 68(3):411--436, 2006.

\bibitem{kong1994} A.~Kong, J.~S. Liu, and W.~H. Wong.
Sequential imputations and Bayesian missing data problems.
\emph{Journal of the American Statistical Association}, 89(425):278--288, 1994.

\bibitem{wilson1927} E.~B. Wilson.
Probable inference, the law of succession, and statistical inference.
\emph{Journal of the American Statistical Association}, 22(158):209--212, 1927.

\bibitem{newcombe1998} R.~G. Newcombe.
Interval estimation for the difference between independent proportions.
\emph{Statistics in Medicine}, 17(8):873--890, 1998.

\bibitem{thompson1933} W.~R. Thompson.
On the likelihood that one unknown probability exceeds another in view of the
evidence of two samples. \emph{Biometrika}, 25(3--4):285--294, 1933.

\bibitem{lattimore2020} T.~Lattimore and C.~Szepesv\'{a}ri.
\emph{Bandit Algorithms}. Cambridge University Press, 2020.

\bibitem{kaelbling1998} L.~P. Kaelbling, M.~L. Littman, and A.~R. Cassandra.
Planning and acting in partially observable stochastic domains.
\emph{Artificial Intelligence}, 101(1--2):99--134, 1998.

\bibitem{lewis2020} P.~Lewis et~al.
Retrieval-augmented generation for knowledge-intensive NLP tasks.
\emph{NeurIPS}, 33:9459--9474, 2020.

\bibitem{yao2023} S.~Yao et~al.
ReAct: Synergizing reasoning and acting in language models. \emph{ICLR}, 2023.

\bibitem{schick2023} T.~Schick et~al.
Toolformer: Language models can teach themselves to use tools.
\emph{NeurIPS}, 36, 2023.

\bibitem{robertson2009} S.~Robertson and H.~Zaragoza.
The probabilistic relevance framework: BM25 and beyond.
\emph{Foundations and Trends in Information Retrieval}, 3(4):333--389, 2009.

\bibitem{deerwester1990} S.~Deerwester et~al.
Indexing by latent semantic analysis.
\emph{Journal of the American Society for Information Science}, 41(6):391--407, 1990.

\end{thebibliography}

\end{document}
